\justifying
\NSFCBodyText

\subsubsection{年度研究计划}
本项目研究期限为2027年1月至2028年12月，按照数据与基线、概率校准、事件检验的依赖关系安排两年任务。下列季度安排为实施草案，具体进度以数据授权、研究合同和人员落实情况调整，\textbf{【暂定首季度完成数据清单与授权核验，第一年完成概率预测基线，第二年完成爬坡评价与原型】}。

\textbf{2027年：完成数据核验、尺度分析与概率预测基线。}项目启动后先核实目标场站、历史记录及气象档案的使用条件，整理时间、功率单位、容量和运行标记，建立按获取时刻截断的数据流程。随后确定训练、验证、校准和测试划分，完成持续性、分位数回归及轻量时序基线，检查日照、天气和非气象异常的分层。具备数据条件后开展逐源输入与尺度消融，再训练有限规模多尺度分位数模型，比较初步校准结果。并核对各天气分组与爬坡事件的独立过程数量及可检验性。年度里程碑为能够复现的样本索引、数据质量报告、基线与尺度增益结果、概率预测及误差轨迹。若年末仍缺少新疆实测，应在中期材料中明确地域验证目标尚未实现，提出基于真实条件的调整方案。

\textbf{2028年：完成天气变化下的校准比较、爬坡检验与成果整理。}在第一年固定的数据规范与基线上，开展多个时间起点的滚动校准，评估标签延迟、天气转换和样本不足的影响。待完整误差轨迹形成后，进行同边际的总体历史秩、天气条件历史秩与独立排序三组比较，检验功率净变化事件概率，再依据验证集选定的告警规则完成独立测试。后期集中整理天气个例、失败案例、敏感性分析与统计区间，构建可回放的离线预警原型，形成论文、技术报告与可移交的程序说明。对未获支持的假设和未完成的数据条件如实记载，不以补选测试时段替代验证。

\subsubsection{预期研究结果}
预期形成一套气象与功率数据处理及评价流程、一套多尺度概率预测与校准实现、一套爬坡风险离线原型及其使用说明。数据成果首先交付字段规范、质量报告和复现实验索引；实测原始数据的公开或移交范围以授权和基金管理要求为准，不承诺未经许可发布场站数据库。方法报告应说明气象输入和尺度选择的有效范围、校准的覆盖与效率、时间依赖在事件预测中的作用，并保存阴性和失败结果。

成果数量按\textbf{【暂定形成研究论文1篇、算法软件或可登记软件成果1项，支持研究生1--2名参与研究】}规划；论文投稿或录用、软件登记和人员培养分别依据真实进度列报，不将预期成果写成已取得成绩。成果署名、第一资助基金和技术档案移交遵照沙漠基金规则落实，依托单位确认事项在其他说明中列明。最终考核以批准合同为准，正文中的设计参数与数量草案不替代正式任务书。
